\section{Subsidy Design under Buyout Exit}\label{sec:subsidy}

This section derives the paper's main normative benchmark. The goal is to isolate the welfare consequences of buyout exit using only the product-market primitives introduced in Section \ref{sec:model} together with the endogenous merger-attempt condition. Planner-specific local and national incidence adjustments are postponed to Section \ref{sec:decentralized}.

\subsection{Primitive benchmark social value}

For project type $j\in\{I,B\}$, define the benchmark value under independent ownership by
\begin{equation}\label{eq:VjD}
    V_j^D \equiv C^D(j)+\pi_I^D(j)+\pi_E^D(j),
\end{equation}
and the benchmark value under approved acquisition by
\begin{equation}\label{eq:VjA}
    V_j^A \equiv C^M(j)+\Pi^M(j).
\end{equation}
Thus the benchmark welfare effect of approved acquisition relative to independent ownership is
\begin{equation}\label{eq:Omegaj}
    \Omega_j \equiv V_j^A-V_j^D.
\end{equation}
Using the definition of private merger surplus,
\begin{equation}\label{eq:Omegajdecomp}
    \Omega_j
    =
    \bigl(C^M(j)-C^D(j)\bigr)+\Delta_j.
\end{equation}
Equation \eqref{eq:Omegajdecomp} makes clear that the benchmark welfare effect of acquisition combines the private merger gain $\Delta_j$ with the product-market consumer-surplus effect.

Let
\begin{equation}
    a_j(q)\equiv \mathbf{1}\{q\Delta_j\ge k\}
\end{equation}
denote the merger-attempt indicator. The primitive-benchmark successful-project value is then
\begin{equation}\label{eq:BjPMq}
    B_j^{PM}(q)
    \equiv
    V_j^D+a_j(q)\bigl[q\Omega_j-k\bigr].
\end{equation}
This object is the benchmark welfare counterpart to the private continuation value in \eqref{eq:Rj}. If no merger is attempted, successful entry yields independent competition and therefore value $V_j^D$. If a merger is attempted, approved acquisition occurs with probability $q$, rejected merger with probability $1-q$, and the resource cost $k$ is incurred regardless of approval.

\subsection{Benchmark welfare conditional on project type}

Fix project type $j$. As in Section \ref{sec:private}, the entrepreneur chooses effort privately, so under Assumption \ref{ass:interior},
\begin{equation}
    x_j^*(q)=R_j(q),
    \qquad
    R_j(q)=\pi_E^D(j)+\beta[q\Delta_j-k]_+.
\end{equation}
For an entrepreneur with entry cost $f$, expected benchmark welfare from entering with project $j$ is therefore
\begin{equation}\label{eq:omegajPM}
    \omega_j^{PM}(f;s,q)
    =
    R_j(q)B_j^{PM}(q)
    -
    \frac{R_j(q)^2}{2}
    -
    f
    -
    \tau s.
\end{equation}
The private entry cutoff remains
\begin{equation}\label{eq:cjqPM}
    c_j(s,q)=s+\frac{R_j(q)^2}{2}.
\end{equation}
Hence benchmark welfare conditional on project $j$ is
\begin{equation}\label{eq:WjPM}
    W_j^{PM}(s,q)
    =
    \int_0^{c_j(s,q)}
    \left[
        R_j(q)B_j^{PM}(q)-\frac{R_j(q)^2}{2}-f-\tau s
    \right]df.
\end{equation}

\begin{proposition}\label{prop:pmbenchmarksubsidy}
Fix a project type $j\in\{I,B\}$ and an approval probability $q$. Then $W_j^{PM}(s,q)$ is strictly concave in $s$, and the unique optimal benchmark subsidy is
\begin{equation}\label{eq:sjPMstar}
    s_j^{PM*}(q)
    =
    \max\left\{
        0,\;
        \frac{
            R_j(q)B_j^{PM}(q)-\frac{\tau+2}{2}R_j(q)^2
        }{2\tau+1}
    \right\}.
\end{equation}
Equivalently,
\begin{equation}
    s_j^{PM*}(q)
    =
    \max\left\{
        0,\;
        \frac{
            R_j(q)\left[B_j^{PM}(q)-\frac{\tau+2}{2}R_j(q)\right]
        }{2\tau+1}
    \right\}.
\end{equation}
\end{proposition}

Proof is reported in Appendix \ref{app:decentralized}.

Proposition \ref{prop:pmbenchmarksubsidy} is the benchmark counterpart to the planner-specific subsidy formulas studied later. Its key feature is that the relevant welfare block is determined entirely by the product-market primitives together with the endogenous merger-attempt condition.

\subsection{Comparative statics under the primitive benchmark}

The first implication is that merger permissiveness matters only once a merger is privately attempted. If $q<\bar q_j\equiv k/\Delta_j$, then $a_j(q)=0$, so
\begin{equation}
    R_j(q)=\pi_E^D(j),
    \qquad
    B_j^{PM}(q)=V_j^D,
\end{equation}
and therefore the benchmark subsidy is locally constant in $q$.

\begin{proposition}\label{prop:pmcomparative}
Fix a project type $j\in\{I,B\}$.
\begin{enumerate}
    \item If $q<\bar q_j\equiv k/\Delta_j$, then $s_j^{PM*}(q)$ is constant in $q$.

    \item Consider an interval on which $q>\bar q_j$ and the optimal subsidy is interior. Then
    \begin{equation}\label{eq:dspmdq}
        \frac{d s_j^{PM*}(q)}{dq}
        =
        \frac{
            \beta\Delta_j\bigl[B_j^{PM}(q)-(\tau+2)R_j(q)\bigr]
            +
            R_j(q)\Omega_j
        }{2\tau+1}.
    \end{equation}

    \item In particular, if
    \begin{equation}\label{eq:pmnegativecase}
        \Omega_j\le 0
        \qquad\text{and}\qquad
        B_j^{PM}(q)\le (\tau+2)R_j(q),
    \end{equation}
    then
    \begin{equation}
        \frac{d s_j^{PM*}(q)}{dq}\le 0.
    \end{equation}
\end{enumerate}
\end{proposition}

Proof is reported in Appendix \ref{app:decentralized}.

Proposition \ref{prop:pmcomparative} isolates the model-driven comparative statics in a transparent way. The relevant objects are the private merger surplus $\Delta_j$, the benchmark social merger gain $\Omega_j$, and the private continuation value $R_j(q)$. No planner-specific incidence parameters appear in the theorem.

\subsection{Uniform subsidy under private project choice}

Return now to the actual equilibrium in which project choice is private and governed by Proposition \ref{prop:projectchoice}. Since the uniform subsidy does not affect project choice, the planner chooses $s$ taking $j^*(q)$ as given. Thus equilibrium benchmark welfare is
\begin{equation}
    W^{PM}(s,q)=W_{j^*(q)}^{PM}(s,q),
\end{equation}
and the optimal uniform benchmark subsidy is
\begin{equation}\label{eq:sPMuniform}
    s^{PM*}(q)=s_{j^*(q)}^{PM*}(q).
\end{equation}

\begin{proposition}\label{prop:pmuniformsubsidy}
Suppose Assumptions \ref{ass:types}--\ref{ass:singlecross} hold. Then the optimal uniform benchmark subsidy is
\begin{equation}
    s^{PM*}(q)
    =
    \begin{cases}
        s_I^{PM*}(q), & q<\hat q,\\
        \{s_I^{PM*}(\hat q),\, s_B^{PM*}(\hat q)\}, & q=\hat q,\\
        s_B^{PM*}(q), & q>\hat q.
    \end{cases}
\end{equation}
Moreover, if
\begin{equation}\label{eq:pmdropzero}
    B_B^{PM}(\hat q)<\frac{\tau+2}{2}R_B(\hat q),
\end{equation}
then there exists $\varepsilon>0$ such that
\begin{equation}
    s^{PM*}(q)=0
    \qquad
    \text{for all } q\in(\hat q,\hat q+\varepsilon).
\end{equation}
\end{proposition}

Proof is reported in Appendix \ref{app:decentralized}.

Proposition \ref{prop:pmuniformsubsidy} gives the paper's benchmark normative result. Merger permissiveness can redirect private project choice toward type $B$, and if the buyout-oriented project has weak benchmark social value relative to its private continuation value, the optimal subsidy falls and may vanish immediately after the project switch.

\subsection{Discussion}

The benchmark results in this section are driven by two primitive objects. The first is $\Delta_j$, which governs private merger feasibility and the private continuation value of entrepreneurship. The second is $\Omega_j=\Delta_j+\bigl(C^M(j)-C^D(j)\bigr)$, which governs the benchmark social gain from approved acquisition. This separation makes clear which comparative statics are model-driven before any planner-specific incidence assumptions are imposed.

Section \ref{sec:decentralized} builds on this benchmark by reintroducing local retention, planner-specific spillovers, and other incidence adjustments. Those planner-specific results therefore modify the benchmark derived here rather than replacing it.
